% figures/w01-pqr.svg — p, q, r 세 각속도.
%
%   bash _tools/tikz2svg.sh figures/src/w01-pqr.tex figures/w01-pqr.svg
%
% 셋 다 선체의 축 하나에 대해 재고, 셋 다 오른손 법칙을 따른다. z 가 아래를
% 향하므로 양의 yaw 는 위에서 보면 시계방향이고, 양의 pitch 는 선수를 든다.
% 각 칸은 그 축을 관측자 쪽으로 두고 본 그림이다.
%
% ── 회전 방향은 계산해서 정한다 (CLAUDE.md §4 규칙 6) ────────────────────────
% 화면 기저를 오른쪽 R, 위 U, 화면 밖 O 로 두고 R x U = O 라 하자.
%
%   (p) 선미에서 앞을 본다.  x_b = -O (화면 안쪽), z_b = -U (아래).
%       x_b x y_b = z_b 이므로 y_b = R — 우현이 화면 오른쪽. 맞다.
%       x_b 둘레의 양의 회전은 y_b 를 z_b 로 보낸다: R -> -U, 즉 오른쪽이 아래로.
%       화면에서 **시계방향**. 호는 위를 지나 왼쪽에서 오른쪽으로 그린다.
%   (q) 우현에서 본다.  y_b = +O (화면 밖).  O 둘레의 양의 회전은 반시계방향.
%       오른쪽(선수)이 올라간다. 호는 위를 지나 오른쪽에서 왼쪽으로.
%   (r) 위에서 본다.  z_b = -O (화면 안쪽).  다시 **시계방향** — 선수가 우현으로.
%
% 출처: SNAME (1950) 표기, Fossen (2021) Handbook §2.1 표 2.1.
% TikZ 에서 arc 의 각이 **줄어들면** 시계방향, 늘어나면 반시계방향이다.
\documentclass[border=5pt]{standalone}
% gnc-style.tex — 이 볼트의 TikZ 그림이 공유하는 스타일.
%
% 저널 그림 규약을 스타일 하나로 굳혀 둔다. 그림마다 선 굵기나 화살촉을
% 다시 정하지 않는다 — 그것이 그림이 제각각으로 보이는 원인이다.
%
% 쓰는 법:  % gnc-style.tex — 이 볼트의 TikZ 그림이 공유하는 스타일.
%
% 저널 그림 규약을 스타일 하나로 굳혀 둔다. 그림마다 선 굵기나 화살촉을
% 다시 정하지 않는다 — 그것이 그림이 제각각으로 보이는 원인이다.
%
% 쓰는 법:  % gnc-style.tex — 이 볼트의 TikZ 그림이 공유하는 스타일.
%
% 저널 그림 규약을 스타일 하나로 굳혀 둔다. 그림마다 선 굵기나 화살촉을
% 다시 정하지 않는다 — 그것이 그림이 제각각으로 보이는 원인이다.
%
% 쓰는 법:  \input{gnc-style}  (figures/src/ 안에서)

\usepackage{tikz}
\usepackage{amsmath}
\usepackage{xcolor}
\usetikzlibrary{arrows.meta, positioning, calc, fit, backgrounds, decorations.pathmorphing}

% ---- 색 --------------------------------------------------------------------
% 색은 강조 하나에만 쓴다. 나머지는 검정.
\definecolor{ink}{HTML}{1A1A1A}
\definecolor{accent}{HTML}{7C3AED}
\definecolor{warn}{HTML}{B91C1C}
\definecolor{soft}{HTML}{666666}

% ---- 선과 화살촉 -----------------------------------------------------------
\tikzset{
  % 신호선. 화살촉은 작고 뾰족한 것 하나뿐이다.
  sig/.style   = {ink, line width=0.5pt, -{Stealth[length=4pt, width=3pt]}},
  wire/.style  = {ink, line width=0.5pt},
  % 강조 경로 — 파선으로 구분한다. 흑백 인쇄에서도 살아야 하기 때문이다.
  asig/.style  = {accent, line width=0.5pt, densely dashed,
                  -{Stealth[length=4pt, width=3pt]}},
  awire/.style = {accent, line width=0.5pt, densely dashed},
  % 블록. 흰 바탕, 직각, 검은 테두리.
  blk/.style   = {draw=ink, line width=0.55pt, fill=white,
                  minimum width=17mm, minimum height=8mm, inner sep=2pt, font=\small},
  % 합산점
  sum/.style   = {draw=ink, line width=0.55pt, fill=white, circle,
                  inner sep=0pt, minimum size=4.6mm, font=\scriptsize},
  asum/.style  = {draw=accent, line width=0.55pt, fill=white, circle,
                  inner sep=0pt, minimum size=4.6mm, font=\scriptsize},
  % 분기점
  dot/.style   = {ink, fill=ink, circle, inner sep=0pt, minimum size=1.5mm},
  % 신호 이름 — 선 위에 작게
  sname/.style = {font=\small, inner sep=1.5pt},
  % 그림 안의 짧은 설명. 그림 안의 글은 최소로 둔다
  note/.style  = {font=\scriptsize, soft, inner sep=1.5pt},
  anote/.style = {font=\scriptsize, accent, inner sep=1.5pt},
  % 축
  axis/.style  = {ink, line width=0.5pt},
  tick/.style  = {font=\scriptsize, soft},
}

% 캡션 한 줄 — 그림 아래 가는 선과 함께
\newcommand{\gnccaption}[2]{%
  \node[anchor=north west, text width=#1, align=left, font=\scriptsize, ink]
        at (current bounding box.south west) {\vspace{2pt}#2};}
  (figures/src/ 안에서)

\usepackage{tikz}
\usepackage{amsmath}
\usepackage{xcolor}
\usetikzlibrary{arrows.meta, positioning, calc, fit, backgrounds, decorations.pathmorphing}

% ---- 색 --------------------------------------------------------------------
% 색은 강조 하나에만 쓴다. 나머지는 검정.
\definecolor{ink}{HTML}{1A1A1A}
\definecolor{accent}{HTML}{7C3AED}
\definecolor{warn}{HTML}{B91C1C}
\definecolor{soft}{HTML}{666666}

% ---- 선과 화살촉 -----------------------------------------------------------
\tikzset{
  % 신호선. 화살촉은 작고 뾰족한 것 하나뿐이다.
  sig/.style   = {ink, line width=0.5pt, -{Stealth[length=4pt, width=3pt]}},
  wire/.style  = {ink, line width=0.5pt},
  % 강조 경로 — 파선으로 구분한다. 흑백 인쇄에서도 살아야 하기 때문이다.
  asig/.style  = {accent, line width=0.5pt, densely dashed,
                  -{Stealth[length=4pt, width=3pt]}},
  awire/.style = {accent, line width=0.5pt, densely dashed},
  % 블록. 흰 바탕, 직각, 검은 테두리.
  blk/.style   = {draw=ink, line width=0.55pt, fill=white,
                  minimum width=17mm, minimum height=8mm, inner sep=2pt, font=\small},
  % 합산점
  sum/.style   = {draw=ink, line width=0.55pt, fill=white, circle,
                  inner sep=0pt, minimum size=4.6mm, font=\scriptsize},
  asum/.style  = {draw=accent, line width=0.55pt, fill=white, circle,
                  inner sep=0pt, minimum size=4.6mm, font=\scriptsize},
  % 분기점
  dot/.style   = {ink, fill=ink, circle, inner sep=0pt, minimum size=1.5mm},
  % 신호 이름 — 선 위에 작게
  sname/.style = {font=\small, inner sep=1.5pt},
  % 그림 안의 짧은 설명. 그림 안의 글은 최소로 둔다
  note/.style  = {font=\scriptsize, soft, inner sep=1.5pt},
  anote/.style = {font=\scriptsize, accent, inner sep=1.5pt},
  % 축
  axis/.style  = {ink, line width=0.5pt},
  tick/.style  = {font=\scriptsize, soft},
}

% 캡션 한 줄 — 그림 아래 가는 선과 함께
\newcommand{\gnccaption}[2]{%
  \node[anchor=north west, text width=#1, align=left, font=\scriptsize, ink]
        at (current bounding box.south west) {\vspace{2pt}#2};}
  (figures/src/ 안에서)

\usepackage{tikz}
\usepackage{amsmath}
\usepackage{xcolor}
\usetikzlibrary{arrows.meta, positioning, calc, fit, backgrounds, decorations.pathmorphing}

% ---- 색 --------------------------------------------------------------------
% 색은 강조 하나에만 쓴다. 나머지는 검정.
\definecolor{ink}{HTML}{1A1A1A}
\definecolor{accent}{HTML}{7C3AED}
\definecolor{warn}{HTML}{B91C1C}
\definecolor{soft}{HTML}{666666}

% ---- 선과 화살촉 -----------------------------------------------------------
\tikzset{
  % 신호선. 화살촉은 작고 뾰족한 것 하나뿐이다.
  sig/.style   = {ink, line width=0.5pt, -{Stealth[length=4pt, width=3pt]}},
  wire/.style  = {ink, line width=0.5pt},
  % 강조 경로 — 파선으로 구분한다. 흑백 인쇄에서도 살아야 하기 때문이다.
  asig/.style  = {accent, line width=0.5pt, densely dashed,
                  -{Stealth[length=4pt, width=3pt]}},
  awire/.style = {accent, line width=0.5pt, densely dashed},
  % 블록. 흰 바탕, 직각, 검은 테두리.
  blk/.style   = {draw=ink, line width=0.55pt, fill=white,
                  minimum width=17mm, minimum height=8mm, inner sep=2pt, font=\small},
  % 합산점
  sum/.style   = {draw=ink, line width=0.55pt, fill=white, circle,
                  inner sep=0pt, minimum size=4.6mm, font=\scriptsize},
  asum/.style  = {draw=accent, line width=0.55pt, fill=white, circle,
                  inner sep=0pt, minimum size=4.6mm, font=\scriptsize},
  % 분기점
  dot/.style   = {ink, fill=ink, circle, inner sep=0pt, minimum size=1.5mm},
  % 신호 이름 — 선 위에 작게
  sname/.style = {font=\small, inner sep=1.5pt},
  % 그림 안의 짧은 설명. 그림 안의 글은 최소로 둔다
  note/.style  = {font=\scriptsize, soft, inner sep=1.5pt},
  anote/.style = {font=\scriptsize, accent, inner sep=1.5pt},
  % 축
  axis/.style  = {ink, line width=0.5pt},
  tick/.style  = {font=\scriptsize, soft},
}

% 캡션 한 줄 — 그림 아래 가는 선과 함께
\newcommand{\gnccaption}[2]{%
  \node[anchor=north west, text width=#1, align=left, font=\scriptsize, ink]
        at (current bounding box.south west) {\vspace{2pt}#2};}


\begin{document}
\begin{tikzpicture}[x=1mm, y=1mm]

\node[ink, font=\small\bfseries, anchor=north west] at (-6,60)
     {$p$, $q$ and $r$ — the three body angular rates};
\node[anchor=north west, align=left, text width=170mm, font=\scriptsize, soft] at (-6,54) {%
  Each is measured about one axis of the \textbf{hull}, by a gyroscope bolted to
  the hull. All three are in rad/s and all three obey the right-hand rule.};

% ===================== p : roll =============================================
\begin{scope}
  \node[ink, font=\small\bfseries, anchor=west] at (-6,41) {$p$ — roll rate};
  \node[note, anchor=west] at (-6,36.5) {seen from astern, looking forward};

  % 선체 단면 : 쌍동선
  \foreach \s in {-1,1} {
    \fill[ink, opacity=0.10] ({\s*11-4},0) rectangle ({\s*11+4},7);
    \draw[ink, line width=0.5pt] ({\s*11-4},0) rectangle ({\s*11+4},7);
  }
  \draw[ink, line width=0.5pt] (-15,7) -- (15,7) -- (15,13) -- (-15,13) -- cycle;

  % x 축은 관측자에서 멀어진다 : 원 안의 x
  \draw[ink, line width=0.5pt] (0,10) circle (2.0);
  \draw[ink, line width=0.5pt] (-1.4,8.6) -- (1.4,11.4);
  \draw[ink, line width=0.5pt] (-1.4,11.4) -- (1.4,8.6);
  \node[font=\scriptsize, anchor=south] at (0,13.8) {$x_b$ — away from you};

  % x 가 화면 안쪽이면 양의 회전은 시계방향 : 각이 줄어드는 호
  \draw[accent, line width=1.0pt, -{Stealth[length=4.5pt,width=3.4pt]}]
        (-19,20) arc (180:10:19 and 7);
  \node[accent, font=\scriptsize, anchor=south] at (0,28.5) {$p > 0$};
  \node[note, anchor=north, align=center, text width=48mm] at (0,-2) {%
    the \textbf{starboard} side goes down;\\ $\dot\phi \approx p$ for a surface craft};
\end{scope}

% ===================== q : pitch ============================================
\begin{scope}[shift={(62,0)}]
  \node[ink, font=\small\bfseries, anchor=west] at (-6,41) {$q$ — pitch rate};
  \node[note, anchor=west] at (-6,36.5) {seen from starboard, bow to the right};

  \fill[ink, opacity=0.10] (-16,5) -- (14,5) -- (19,11) -- (-16,11) -- cycle;
  \draw[ink, line width=0.5pt] (-16,5) -- (14,5) -- (19,11) -- (-16,11) -- cycle;

  % y 축은 관측자 쪽 : 원 안의 점
  \draw[ink, line width=0.5pt] (0,8) circle (2.0);
  \fill[ink] (0,8) circle (0.7);
  \node[font=\scriptsize, anchor=south] at (0,11.8) {$y_b$ — towards you};

  % y 가 화면 밖이면 양의 회전은 반시계방향 : 각이 늘어나는 호
  \draw[accent, line width=1.0pt, -{Stealth[length=4.5pt,width=3.4pt]}]
        (19,20) arc (0:170:19 and 7);
  \node[accent, font=\scriptsize, anchor=south] at (0,28.5) {$q > 0$};
  \node[note, anchor=north, align=center, text width=48mm] at (0,-2) {%
    the bow goes \textbf{up} —\\ even though $z$ points down};
\end{scope}

% ===================== r : yaw ==============================================
\begin{scope}[shift={(126,0)}]
  \node[ink, font=\small\bfseries, anchor=west] at (-6,41) {$r$ — yaw rate};
  \node[note, anchor=west] at (-6,36.5) {seen from above, bow upwards};

  \fill[ink, opacity=0.10] (0,18) -- (-7,4) -- (-7,-6) -- (7,-6) -- (7,4) -- cycle;
  \draw[ink, line width=0.5pt] (0,18) -- (-7,4) -- (-7,-6) -- (7,-6) -- (7,4) -- cycle;

  % z 축은 화면 안쪽 (아래) : 원 안의 x
  \draw[ink, line width=0.5pt] (0,7) circle (2.0);
  \draw[ink, line width=0.5pt] (-1.4,5.6) -- (1.4,8.4);
  \draw[ink, line width=0.5pt] (-1.4,8.4) -- (1.4,5.6);
  \node[font=\scriptsize, anchor=west] at (6.8,7) {$z_b$ — into the page};

  % z 가 화면 안쪽이면 양의 회전은 시계방향 : 각이 줄어드는 호
  \draw[accent, line width=1.0pt, -{Stealth[length=4.5pt,width=3.4pt]}]
        (-19,20) arc (180:10:19 and 7);
  \node[accent, font=\scriptsize, anchor=south] at (0,28.5) {$r > 0$};
  \node[note, anchor=north, align=center, text width=48mm] at (0,-8) {%
    the bow swings to \textbf{starboard};\\ $\dot\psi \approx r$ for a surface craft};
\end{scope}

% ---- 아래 한 줄 -------------------------------------------------------------
\draw[soft, line width=0.3pt] (-6,-22) -- (166,-22);
\node[anchor=north west, align=left, text width=172mm, font=\scriptsize] at (-6,-24) {%
  Point the right thumb along the positive axis; the fingers curl in the positive
  direction of rotation. With $z$ \textbf{down}, that makes a positive roll and a
  positive yaw look \textbf{clockwise} on these two views, while a positive pitch
  looks anticlockwise from starboard and lifts the bow. Both feel wrong once and
  then never again. The convention is SNAME (1950); Fossen (2021) Handbook §2.1,
  Table 2.1 states it in this form.\\[3pt]
  These are body rates, not Euler-angle rates: §1-5 gives the matrix
  $\mathbf{T}_\Theta$ that separates them, and it is \textbf{not} a rotation.};

\end{tikzpicture}
\end{document}
